% --- 2.1.tex ---
\subsection{Nhận dạng thực thể tên riêng (Named Entity Recognition - NER)}

\subsubsection{Mục tiêu bài toán}
Tác vụ này hướng tới việc tự động nhận diện và phân loại các thực thể quan trọng trong văn bản hợp đồng vào các nhóm định trước như: Tên tổ chức/cá nhân (Bên A, Bên B), Ngày tháng, Địa điểm, Giá trị hợp đồng, v.v. Đây là bước then chốt để chuyển đổi văn bản phi cấu trúc thành thông tin có cấu trúc, phục vụ cho việc quản lý và tra cứu hợp đồng tự động.

\subsubsection{Phương pháp tiếp cận: Mô hình PhoBERT}
Nhóm lựa chọn sử dụng \textbf{PhoBERT} (\texttt{vinai/phobert-base}), một mô hình Transformer được huấn luyện trước (pre-trained) dành riêng cho tiếng Việt, dựa trên kiến trúc RoBERTa. PhoBERT cho thấy hiệu năng vượt trội trong các tác vụ NLP tiếng Việt nhờ khả năng hiểu ngữ cảnh và đặc thù từ ghép.

\subsubsection{Quy trình triển khai}
Quá trình xây dựng mô hình được thực hiện qua các giai đoạn sau:

\begin{enumerate}
    \item \textbf{Chuẩn bị dữ liệu và Gán nhãn (BIO Tagging):}
    Dữ liệu được chuẩn hóa về định dạng BIO (Begin, Inside, Outside). Mỗi token trong câu được gán một nhãn tương ứng, ví dụ: \texttt{B-ORG} cho từ bắt đầu tên tổ chức, \texttt{I-ORG} cho các từ tiếp theo và \texttt{O} cho các từ không thuộc thực thể nào.

    \item \textbf{Tiền xử lý (Preprocessing):}
    \begin{itemize}
        \item \textbf{Phân từ:} Sử dụng VnCoreNLP để phân từ trước khi đưa vào mô hình để đảm bảo tính nhất quán với quá trình pre-train của PhoBERT.
        \item \textbf{Tokenization:} Sử dụng \texttt{AutoTokenizer} của thư viện \texttt{transformers} để chuyển đổi các câu thành các ID vector, đồng thời xử lý vấn đề \textit{subword} bằng cách căn chỉnh nhãn (label alignment) cho các token bị chia nhỏ.
    \end{itemize}

    \item \textbf{Kiến trúc mô hình và Huấn luyện:}
    \begin{itemize}
        \item Nhóm sử dụng lớp \texttt{AutoModelForTokenClassification} để tích hợp một lớp tuyến tính (Linear layer) lên trên kiến trúc PhoBERT nhằm thực hiện phân loại thực thể cho từng token.
        \item \textbf{Tham số huấn luyện:} Sử dụng bộ tối ưu hóa AdamW, tốc độ học (learning rate) thấp (ví dụ: $5 \times 10^{-5}$) và kỹ thuật \textit{Early Stopping} để tránh hiện tượng quá khớp (overfitting).
    \end{itemize}

    \item \textbf{Đánh giá:}
    Mô hình được đánh giá dựa trên các chỉ số \textit{Precision}, \textit{Recall} và \textit{F1-score} trên tập dữ liệu kiểm tra (test set).
\end{enumerate}

\subsubsection{Kết quả đạt được}
Mô hình PhoBERT đạt được độ chính xác cao trong việc nhận diện các thực thể có cấu trúc ổn định như Ngày tháng và Giá trị tiền tệ. Đối với các thực thể phức tạp như Tên công ty hoặc Nội dung điều khoản, mô hình vẫn thể hiện khả năng trích xuất tốt nhờ vào việc hiểu ngữ cảnh xung quanh thực thể trong mệnh đề pháp lý.